\section{Training, reuse, and policy-cost precision}\label{sec:study70}
The experiment links the constructive readout to the actual controlled economy. It compares representation and training loss at a common information budget, validates a frozen readout on independent contexts, measures a longer deployment window, and evaluates the returned controllers at a prespecified policy-cost margin. The original continuous laws, state-dependent investment capacity and Bellman target are unchanged.

\subsection{Design and information balance}
The scalar tasks are $(d,T,\varepsilon)=(2,2,1/128)$ and $(8,3,1/64)$, with twenty-four-bit state acquisition and thirty-two-bit downward action rounding. Four fitted proposals receive the same forty-eight training states, approximate action labels and verified boundary endpoints per nonterminal date. The certificate-aware hat--ReLU and degree-four Bernstein readouts have exactly $5d$ box-constrained coefficients and the same objective, penalty and optimizer budget. An action-label readout on the same hat features separates the objective from the representation. A quadratic label fit retains its explicitly different parameter count. Adaptive parabolic search and upper-chord minimization are prediction-free controls; unused training is not charged to them.

The hat basis has an exact finite ReLU realization. It also has a conventional piecewise-linear interpretation, so this comparison does not count two encodings of that basis as independent algorithms. The Bernstein comparison holds the coefficient count and purchased information fixed. All fitted proposals use the same final certificate. Every returned certificate-aware iterate is quantized and receives an exact rational value of the objective and box dual gap from Proposition~\ref{prop:readout67}. The achieved gap is reported rather than assuming that an optimizer's success flag proves optimality.

The protocol fixes two fitting seeds and four process repetitions. It contains eighty cold comparisons, sixteen independent-validation services, fifty-six reuse services, two paired-inference services and thirty-four vector services: 188 complete processes. All returned. This is finite-catalogue evidence for the complete verifier, not a population claim about fitting success. Four clock repetitions are not four new training samples. The construction, receipts and original source freeze predate this manuscript integration; the publication replay adds no scientific observations.

\begin{table}[htbp]
\centering\caption{Complete cold-start services at a common information budget}\label{tab:cold70}
{\small\setlength{\tabcolsep}{3pt}
\begin{tabular}{lcrrcc}
\hline
$(d,T)$ & Rule & Processes & Policy queries & Median (s) & Range (s) \\
\hline
2,2 & A & 4 & 597 & 0.365 & [0.364, 0.366] \\
2,2 & B & 8 & 607--612 & 0.465 & [0.465, 0.466] \\
2,2 & H & 8 & 592--612 & 0.465 & [0.465, 0.466] \\
2,2 & L & 8 & 592--597 & 0.440 & [0.414, 0.466] \\
2,2 & Q & 8 & 597 & 0.415 & [0.415, 0.967] \\
2,2 & U & 4 & 597 & 0.390 & [0.364, 0.415] \\
8,3 & A & 4 & 2594 & 0.666 & [0.615, 0.666] \\
8,3 & B & 8 & 2594 & 1.618 & [1.618, 1.619] \\
8,3 & H & 8 & 2594 & 1.618 & [1.618, 1.669] \\
8,3 & L & 8 & 2594 & 1.518 & [1.517, 1.568] \\
8,3 & Q & 8 & 2594 & 1.518 & [1.467, 1.519] \\
8,3 & U & 4 & 2594 & 0.641 & [0.615, 0.666] \\
\hline\end{tabular}}
\par\smallskip{\footnotesize A: adaptive; U: upper-chord minimizer; H: certificate-aware hat--ReLU; B: certificate-aware Bernstein; L: label-loss hat--ReLU; Q: quadratic. Queries concern the policy workload; training and validation queries are additional. Four process repetitions per seed are descriptive clocks, not independent fits or confidence bounds. All complete times include launch and shutdown.}
\end{table}

At a fixed purchased-information and coefficient budget, the hat--ReLU readouts close 48 and 50 of 256 small-task validation contexts for the two seeds, compared with 25 and 21 for Bernstein. The achieved exact box dual gaps are at most $10^{-7}$ across the certificate-aware catalogue. This is a method-specific improvement in the declared context screen, not a uniform complete-runtime advantage. All methods and both tasks remain in the supplementary validation table.


Cold-start work remains an exacting comparison. In the two-state task the median complete hat--ReLU process takes about $0.465$ seconds, against $0.365$ for adaptive search. In the eight-state task the respective medians are about $1.618$ and $0.666$ seconds; both use 2,594 policy-query evaluations in the displayed workload. These counts exclude separately reported training queries. Thus the certified objective is now an executed arm, but its execution is not an identified complete-work improvement. The recorded fitting, prediction and descendant counts locate the cost rather than attributing every successful return to a successful predictor.

\subsection{Prospective validation and reuse}
Each fitted readout is frozen before a new predictor-independent context stream drawn uniformly from the sixteen-bit dyadic state grid is generated. There are 256 contexts per nonterminal date. The validation calculation executes the declared complete-action mesh recovery whenever the pre-query screen fails. Its clipped loss bounds the additional root queries on every recorded context. The supplementary table gives the achieved loss, actual recovery queries, simultaneous validation upper bound and descendant counts for every task, date, method and seed. These 6,144 method--context--date evaluations validate the stated context law and fallback; they are not relabeled as independent on-policy recursive states.

Each reuse process trains and serializes once and then serves sixty-four predetermined blocks of eight trajectories. At block index 16 it deliberately receives a stale acquisition identifier and uses the adaptive benchmark. Loading, hashing, identity validation, rejection, queries, cash-path enclosures, durable files and process shutdown are included. A stable observed crossing must begin by block index 48 and remain below the matched adaptive prefix through index 63. This requirement is fixed before execution.

\begin{table}[htbp]
\centering\caption{Sixty-four-block prospective reuse}\label{tab:reuse70}
{\small\setlength{\tabcolsep}{3pt}
\begin{tabular}{lcrcccc}
\hline
$(d,T)$ & Rule & Policy queries & Median (s) & Range (s) & Wins & Stable \\
\hline
2,2 & A & 10057 & 3.621 & [3.169, 7.129] & -- & -- \\
2,2 & B & 10392--10502 & 3.521 & [3.320, 3.572] & 4/8 & 4/8 \\
2,2 & H & 10332--10447 & 3.746 & [3.320, 5.075] & 4/8 & 3/8 \\
2,2 & Q & 9982--9992 & 3.621 & [3.120, 4.473] & 4/8 & 3/8 \\
8,3 & A & 218656 & 24.914 & [24.762, 25.216] & -- & -- \\
8,3 & B & 220282--220579 & 26.444 & [26.068, 26.921] & 0/8 & 0/8 \\
8,3 & H & 220498--222019 & 26.469 & [25.915, 27.568] & 0/8 & 0/8 \\
8,3 & Q & 223929--226193 & 27.872 & [26.964, 28.973] & 0/8 & 0/8 \\
\hline\end{tabular}}
\par\smallskip{\footnotesize A complete-time win is relative to the adaptive process with the same repetition index. Stable means an observed cumulative crossing by index 48 that persists through index 63. The last summary and shutdown belong to complete clocks, not prefix clocks. The stale identity at index 16 is declined in every process. The same adaptive repetition is shared across matched seeds; these comparisons are not independent trials.}
\end{table}


Three of the eight hat--ReLU processes in the two-state task satisfy that stable-prefix criterion; four have smaller final complete times than their matched adaptive process. The Bernstein and quadratic rules also exhibit small-task crossings. None of the learned processes wins the eight-state reuse comparison. In the small task, the adaptive complete times range from about $3.17$ to $7.13$ seconds, while hat--ReLU ranges from about $3.32$ to $5.07$ seconds. Consequently, a favorable matched crossing cannot be interpreted as uniform amortized superiority. The operation counts and every unfavorable repetition remain in the record.

The deployment theorem states what additional context-law and descendant-cost information would be sufficient for a prospective work bound. The experiment does not substitute its identity check for a distributional drift bound. Likewise, the sixty-four-block window is not extrapolated to an unobserved stopping horizon.

\subsection{A cost comparison at the declared precision}
The central inference change is the Bellman-centered observation of Theorem~\ref{thm:centered69}. It estimates the actual expected cost difference by evaluating optimal advantages along each actual policy trajectory. The conditional Bellman expectations are evaluated with a tighter tolerance and paid for separately. The estimator is not the difference between two policy-loss upper bounds, and its realized observation is not asserted equal to a same-shock cash-cost difference.

For each scalar task the protocol fixes 1,024 new common-path rows, forty-eight-bit continuous bins and evaluation tolerance $\varepsilon/32$. The policies are adaptive search, certificate-aware hat--ReLU, certificate-aware Bernstein, and the original R67 ReLU controller at its pinned fitting seed. That historical controller is loaded unchanged. The predeclared loss margin is $1/256$, with numerical and sampling half-width allocations of $1/512$ each. No new row is added after inspecting a sign or a failed comparison. There are twelve paired contrasts across the two tasks on one centered-inference event with failure probability at most $0.01$.

\begin{table}[htbp]
\centering\caption{Bellman-centered expected policy-cost differences}\label{tab:inference70}
{\small\setlength{\tabcolsep}{3pt}
\begin{tabular}{lccrrc}
\hline
$(d,T)$ & Pair & Simultaneous interval & Numerical width & Radius & Margin \\
\hline
2,2 & A--B & [-0.000683, 0.000709] & 0.000194 & 0.000599 & Yes \\
2,2 & A--H & [-0.000683, 0.000708] & 0.000194 & 0.000599 & Yes \\
2,2 & A--R & [-0.000658, 0.000728] & 0.000188 & 0.000599 & Yes \\
2,2 & B--R & [-0.000670, 0.000715] & 0.000187 & 0.000599 & Yes \\
2,2 & H--B & [-0.000695, 0.000695] & 0.000193 & 0.000599 & Yes \\
2,2 & H--R & [-0.000670, 0.000715] & 0.000187 & 0.000599 & Yes \\
8,3 & A--B & [-0.001603, 0.001603] & 0.000812 & 0.001197 & Yes \\
8,3 & A--H & [-0.001603, 0.001603] & 0.000812 & 0.001197 & Yes \\
8,3 & A--R & [-0.001583, 0.001613] & 0.000801 & 0.001197 & Yes \\
8,3 & B--R & [-0.001583, 0.001613] & 0.000802 & 0.001197 & Yes \\
8,3 & H--B & [-0.001603, 0.001603] & 0.000812 & 0.001197 & Yes \\
8,3 & H--R & [-0.001583, 0.001613] & 0.000802 & 0.001197 & Yes \\
\hline\end{tabular}}
\par\smallskip{\footnotesize Pair first minus second; R is the unchanged R67 ReLU readout at seed 6603. Every task uses 1,024 new common-path rows, not that number per pair. All intervals are outward-rounded. Margin denotes containment in the prospectively specified interval $[-1/256,1/256]$, not strict superiority. The numerical term is a full width; the sampling term is a half-width.}
\end{table}


All twelve centered intervals lie within $[-1/256,1/256]$. They therefore establish two-sided equivalence at the declared normalized-loss margin, including comparisons involving the unchanged historical neural controller. Every interval still contains zero: strict cost superiority is not identified. Both precision allocations are met. This replaces the earlier inability to assess the declared margin by a positive, original-law eligibility result, without changing the economic estimand or assuming a fitted value is exact.

For adaptive-minus-hat--ReLU, the two-state interval is approximately $[-0.00068263,0.00070764]$ and the eight-state interval is approximately $[-0.00160291,0.00160253]$, with outward display rounding. The corresponding cash-cost intervals on the same new rows have widths about $0.423$ and $0.529$. The centered widths are about $0.00139$ and $0.00321$. This is an identified precision improvement for the same expected-cost target, not an identified neural speed advantage. The complete shared inference processes cost about $9.03$ and $408.63$ seconds, respectively. Conditional Bellman integration, rather than state-bin width, dominates their numerical enclosure budget.

The two controllers' common initial-value terms cancel in expectation. All compatible true states within a sampled bin remain covered. No new inference trajectory in this catalogue has an unresolved acquisition branch. The raw conditional-value width, propagated bin width, acquisition allowance and final arithmetic diagnostics are reported separately; they are not assumed statistically independent, and pre-intersection component widths are not falsely presented as an additive decomposition of the tighter final interval.

\subsection{Two controls and the economic service}
The vector study compares simplicial search, largest-edge uniform refinement, a trained two-output hat readout and a trained two-output quadratic proposal at two state dimensions and horizons two and three. Both learned rules purchase the same twelve teacher states per nonterminal date. Their training objective is action-label loss, not the scalar convex certificate loss. All receive the same triangular lower certificate and feasible rounding rule under two independent continuous innovations. The horizon-four hat proposal transfers its three-date model and is identified as a transfer, not as a four-date fitted model.

\begin{table}[htbp]
\centering\caption{Two-control and two-innovation complete services}\label{tab:vector70}
{\small\setlength{\tabcolsep}{3pt}
\begin{tabular}{lcrrrr}
\hline
$(d,T)$ & Rule & Processes & Queries & Median (s) & Peak KiB \\
\hline
2,2 & H & 2 & 98 & 0.47 & 65520 \\
2,2 & Q & 2 & 99 & 0.49 & 65524 \\
2,2 & S & 2 & 82 & 0.39 & 65164 \\
2,2 & U & 2 & 82 & 0.36 & 65328 \\
2,3 & H & 2 & 5634 & 5.83 & 66004 \\
2,3 & Q & 2 & 5914 & 5.93 & 66256 \\
2,3 & S & 2 & 4541 & 3.70 & 65184 \\
2,3 & U & 2 & 9039 & 7.93 & 65016 \\
2,4 & H & 1 & 872033 & 62.01 & 65872 \\
2,4 & S & 1 & 592658 & 41.37 & 65256 \\
8,2 & H & 2 & 91 & 0.46 & 65780 \\
8,2 & Q & 2 & 91 & 0.46 & 65864 \\
8,2 & S & 2 & 75 & 0.36 & 65252 \\
8,2 & U & 2 & 75 & 0.36 & 65336 \\
8,3 & H & 2 & 3495 & 5.93 & 66028 \\
8,3 & Q & 2 & 3495 & 5.85 & 66208 \\
8,3 & S & 2 & 2645 & 3.55 & 65168 \\
8,3 & U & 2 & 3278 & 4.60 & 65544 \\
\hline\end{tabular}}
\par\smallskip{\footnotesize S: simplicial; U: largest-edge uniform refinement; H and Q: trained two-output hat and quadratic proposals. Ordinary services have eight continuous-bin paths. Horizon-four services have one path, and H transfers its three-date readout. These are workload records, not precision-designed estimates of population policy costs. Peak memory is the maximum observed resident set within the listed processes.}
\end{table}


The vector records contain actual continuous-bin policy paths, complete clocks, peak resident memory and descendant counts. Eight paths per ordinary service, and one per horizon-four transfer service, are workload diagnostics rather than a precision-designed population comparison. The experiment adds a trained vector comparison, but it does not identify a vector expected-cost ranking or remove continuation-tree growth. Derivative branch-and-bound is not an executed arm of this catalogue.

For the scalar deployment decision, the release verifies that the cost-inference readouts coincide with the relevant cold and reuse readouts at seed 6901. Parameter matching is necessary for pairing these records. Under the additional same-numerical-policy premise, Proposition~\ref{prop:adoption70} gives their complete-service adoption region, with normal-policy multiplier 504 and observed work differences. Cross-platform prediction portability is separately discussed in the supplement. The supplementary table reports the work savings and critical price ratios, including cases in which savings are negative. The loss margin is a prospective numerical design choice, not an estimated installation charge. Equivalence at that margin makes policy quality an eligible comparison; actual adoption still depends on complete work and the purchaser's declared prices.

\subsection{Earlier evidence and independent reconstruction}
The original R67 catalogue remains a separate, unchanged companion, together with all preceding adverse studies. It records 212 returns, no neural complete-reuse winner, unresolved ordinary cost intervals, forty-eight neural fitting warnings and a severe horizon-four vector expansion. Those observations are not overwritten by the new loss, new inference rows or different reuse window. The previous accumulation-aliasing counterexample and corrected execution also remain unchanged.

The new publication audit verifies every R69 receipt and stored-file identity, recomputes the achieved rational training certificates, checks every validation loss and executes independently written rational path calculations. A second interval implementation uses exact full-bin state and cost bounds, not only representative points, for unambiguous acquired paths. Its coverage and every diagnostic disagreement are recorded explicitly. This provides independent reconstruction of the sampled cash-path calculation and independent final inference arithmetic. It is not represented as an independent generation of every optimal-advantage endpoint or of the largest vector Bellman descendant tree.
